% ============================================================
% HSBI Beamer — ISLP Lecture Series
% Mock Exam 1 — Solutions (Lectures 1-4, Chapters 1-3)
% Compile: pdflatex -jobname=Mock_Exam_1_Solutions_Slides solutions_slides_1.tex
% ============================================================
\documentclass[aspectratio=169,11pt]{beamer}
\usepackage[utf8]{inputenc}\usepackage[T1]{fontenc}\usepackage[english]{babel}
\usepackage{graphicx}\usepackage{booktabs}\usepackage{amsmath,amssymb,amsfonts}
\usepackage{mathtools}\usepackage{hyperref}\usepackage{xcolor}
\usepackage{verbatim}\usepackage{alltt}
\usepackage{tikz}\usetikzlibrary{calc,arrows.meta,positioning,shapes}
\usepackage{tcolorbox}\tcbuselibrary{skins}

\usetheme{Madrid}\usecolortheme{default}
\setbeamertemplate{navigation symbols}{}
\setbeamertemplate{footline}[frame number]
\setbeamertemplate{blocks}[rounded][shadow=false]
\setbeamertemplate{headline}{%
  \leavevmode\hbox{%
    \begin{beamercolorbox}[wd=\paperwidth,ht=2.2ex,dp=0.9ex]{section in head/foot}%
      {\fontsize{4.6}{5.2}\selectfont%
       \insertsectionnavigationhorizontal{\paperwidth}{}{\hskip0pt plus1filll}}%
    \end{beamercolorbox}}%
}
\setbeamerfont{section in head/foot}{size=\tiny}
\setbeamerfont{palette primary}{size=\tiny}
\setbeamerfont{palette tertiary}{size=\tiny}
\setbeamerfont{section in toc}{size=\footnotesize}

\usepackage{listings}
\definecolor{pyKw}{RGB}{0,0,180}\definecolor{pyStr}{RGB}{20,120,40}
\definecolor{pyCom}{RGB}{120,120,120}\definecolor{accent}{RGB}{38,70,140}
\lstdefinestyle{Pystyle}{language=Python,basicstyle=\ttfamily\footnotesize,
  keywordstyle=\color{pyKw}\bfseries,commentstyle=\color{pyCom}\itshape,
  stringstyle=\color{pyStr},showstringspaces=false,upquote=true,
  columns=fullflexible,breaklines=true,literate={~}{{$\sim$}}1}
\lstset{style=Pystyle}

\newtcolorbox{takeaway}[1][Takeaway]{enhanced,colback=green!4,colframe=green!50!black,
  boxrule=0pt,leftrule=2.5pt,arc=2pt,fonttitle=\bfseries,title=#1,
  left=2mm,right=2mm,top=1mm,bottom=1mm}
\newtcolorbox{numexample}[1][Worked example]{enhanced,colback=orange!4,colframe=orange!70!black,
  boxrule=0pt,leftrule=2.5pt,arc=2pt,fonttitle=\bfseries,title=#1,
  left=2mm,right=2mm,top=1mm,bottom=1mm}
\newtcolorbox{readme}[1][How to read this]{enhanced,colback=blue!4,colframe=accent,
  boxrule=0pt,leftrule=2.5pt,arc=2pt,fonttitle=\bfseries,title=#1,
  left=2mm,right=2mm,top=1mm,bottom=1mm}

\title{Mock Exam 1 --- Solutions}
\subtitle{Lectures 1--4 (Chapters 1--3)}
\author{Prof.\ Dr.\ Christoph Weisser}\institute{Quantitative Research Methods \textperiodcentered\ HSBI}
\date{Summer Semester 2026}\graphicspath{{images/}}

\newtcolorbox{exercise}[1][Exercise]{enhanced, fontupper=\footnotesize, colback=purple!4, colframe=purple!55!black, boxrule=0pt, leftrule=2.5pt, arc=2pt, fonttitle=\bfseries, title=#1, left=2mm, right=2mm, top=1mm, bottom=1mm}
\newtcolorbox{solutionbox}[1][Solution]{enhanced, fontupper=\footnotesize, colback=teal!5, colframe=teal!55!black, boxrule=0pt, leftrule=2.5pt, arc=2pt, fonttitle=\bfseries, title=#1, left=2mm, right=2mm, top=1mm, bottom=1mm}

% Reclaim ~2pt of body height (custom nav headline makes full slides tip over by ~1.83pt)
\addtobeamertemplate{frametitle}{}{\vspace{-2pt}}

\begin{document}
\begin{frame}[plain]\titlepage\end{frame}

% ------------------------------------------------------------
\begin{frame}{Overview of the mock exam}
  \footnotesize
  \begin{center}
  \begin{tabular}{@{}clc@{}}
    \toprule
    \textbf{Problem} & \textbf{Topic} & \textbf{Points}\\
    \midrule
    P1 & Fundamentals of statistical learning & 15\\
    P2 & Bias--variance decomposition & 15\\
    P3 & $K$-nearest neighbours classification & 12\\
    P4 & Simple linear regression by hand & 24\\
    P5 & Multiple regression --- interpreting Python output & 24\\
    \midrule
    $\Sigma$ & & \textbf{90}\\
    \bottomrule
  \end{tabular}
  \end{center}
  \vspace{2pt}
  \begin{takeaway}[Exam conditions]
    90 minutes for 90 points ($\approx$ 1 point per minute). Allowed aids: non-programmable
    calculator and one handwritten A4 sheet (both sides). Justify every answer; round final
    results to 3 decimals.
  \end{takeaway}
\end{frame}

% ============================================================
% P1
% ============================================================
\begin{frame}{Problem 1 --- task}
  \begin{exercise}[Problem 1 (15 P) --- Fundamentals of statistical learning]
    \begin{itemize}
      \item[(a)] \textbf{[4 P]} Prediction vs.\ inference --- explain the difference; one example each.
      \item[(b)] \textbf{[4 P]} Parametric vs.\ non-parametric estimation of $f$; one advantage of each;
        which flexibility--interpretability trade-off is behind the choice?
      \item[(c)] \textbf{[3 P]} Classify: (i) churn prediction for 800 telecom customers;
        (ii) which firm characteristics are associated with CEO salary (90 firms);
        (iii) find groups among 5{,}000 shopping baskets.
      \item[(d)] \textbf{[4 P]} Why is training MSE generally smaller than test MSE? What is overfitting?
    \end{itemize}
  \end{exercise}
\end{frame}

\begin{frame}{Problem 1 --- solution (a): prediction vs.\ inference}
  \begin{solutionbox}[Solution P1(a) --- prediction vs.\ inference (4 P)]
    We estimate $f$ in $Y=f(X)+\varepsilon$; the \emph{goal} decides what we need from $\hat f$.
    \medskip

    \textbf{Prediction.} We only want an \emph{accurate output} $\hat Y=\hat f(X)$ for new inputs
    $X$. The form of $\hat f$ is secondary --- it may be treated as a \emph{black box}, because
    the accuracy of $\hat Y$ is all that matters.\\[0.5mm]
    \emph{Example:} Is an incoming e-mail spam or not? (Or: predict next month's product demand.)
    \medskip

    \textbf{Inference.} We want to understand \emph{how} $Y$ changes with $X_1,\dots,X_p$:
    \emph{which} predictors matter, in \emph{which direction}, and \emph{how strongly}. Here
    $\hat f$ must be \emph{interpretable}.\\[0.5mm]
    \emph{Example:} Which advertising channels increase sales, and by how much per euro spent?
  \end{solutionbox}
  \begin{takeaway}[In one line]
    Prediction targets the \emph{output} $\hat Y$; inference targets the \emph{mechanism} inside $\hat f$.
  \end{takeaway}
\end{frame}

\begin{frame}{Problem 1 --- solution (b): parametric vs.\ non-parametric}
  \begin{solutionbox}[Solution P1(b) --- two ways to estimate $f$ (4 P)]
    \textbf{Parametric --- a two-step approach.}
    \begin{itemize}
      \item[(1)] \emph{Assume a functional form} for $f$, e.g.\ the linear model
        $f(X)=\beta_0+\beta_1X_1+\dots+\beta_pX_p$.
      \item[(2)] \emph{Estimate} the finite set of parameters (e.g.\ by least squares).
    \end{itemize}
    Estimating a whole function is thus reduced to estimating a few numbers.
    \emph{Advantage:} needs relatively few observations; simple to fit and easy to interpret.
    \medskip

    \textbf{Non-parametric.} No explicit form is assumed; $\hat f$ is built to follow the data
    directly (e.g.\ $K$-nearest neighbours, splines).
    \emph{Advantage:} captures a far wider range of shapes of $f$ --- no risk of imposing a badly
    wrong form.
    \medskip

    \textbf{Flexibility--interpretability trade-off.} More flexible methods can reduce bias, but
    need many observations, risk overfitting and are harder to interpret; when \emph{inference}
    is the goal we usually prefer the more restrictive, interpretable models.
  \end{solutionbox}
\end{frame}

\begin{frame}{Problem 1 --- solution (c): classifying the scenarios}
  \begin{solutionbox}[Solution P1(c) --- regression / classification / unsupervised (3 P)]
    \begin{itemize}
      \item[(i)] \textbf{Classification} (supervised). The response ``cancels: yes/no'' is
        \emph{qualitative}; the goal is prediction.
      \item[(ii)] \textbf{Regression} (supervised). The response CEO salary is \emph{quantitative};
        the goal is inference --- which factors are associated with salary, and how strongly.
      \item[(iii)] \textbf{Unsupervised learning}. There is \emph{no response variable}; the goal
        is to find clusters of similar customers.
    \end{itemize}
  \end{solutionbox}
  \begin{takeaway}[Decision rule]
    Is there a response? \emph{Qualitative} $\to$ classification; \emph{quantitative} $\to$
    regression; \emph{none} $\to$ unsupervised.
  \end{takeaway}
\end{frame}

\begin{frame}{Problem 1 --- solution (d): training vs.\ test MSE, overfitting}
  \begin{solutionbox}[Solution P1(d) --- why training MSE $<$ test MSE (4 P)]
    \textbf{Why the training MSE is smaller.} The method is fitted by making the \emph{training}
    error small, so $\hat f$ adapts to the training observations --- \emph{including their random
    noise}. The training MSE therefore \emph{systematically underestimates} the error on fresh
    data, and for a given data set training MSE $\le$ test MSE in general.
    \medskip

    \textbf{Behaviour vs.\ flexibility.} As flexibility grows, the training MSE decreases
    \emph{monotonically}, while the test MSE first decreases and then increases (U-shape).
    \medskip

    \textbf{Overfitting.} The method follows patterns in the \emph{training noise} that do not
    exist in the population; a \emph{less} flexible method would then achieve a \emph{smaller
    test} MSE. Symptom: a widening gap --- very small training MSE together with a rising test MSE.
  \end{solutionbox}
\end{frame}

% ============================================================
% P2
% ============================================================
\begin{frame}{Problem 2 --- task}
  \begin{exercise}[Problem 2 (15 P) --- Bias--variance decomposition]
    At a fixed $x_0$ three methods were analysed; $\sigma^2=\operatorname{Var}(\varepsilon)=1.00$:
    \begin{center}\scriptsize
    \begin{tabular}{lccc}
      \toprule
      & A (linear) & B (moderate) & C (very flexible)\\
      \midrule
      $[\operatorname{Bias}]^2$ & 2.25 & 0.16 & 0.01\\
      $\operatorname{Var}(\hat f(x_0))$ & 0.25 & 1.44 & 3.60\\
      \bottomrule
    \end{tabular}
    \end{center}
    (a) \textbf{[4 P]} State the decomposition of the expected test MSE and explain the components.
    (b) \textbf{[5 P]} Compute the expected test MSE for A--C; which method to choose?
    (c) \textbf{[3 P]} Smallest expected test MSE any method could achieve here?
    (d) \textbf{[3 P]} Behaviour of bias$^2$/variance/training/test MSE as flexibility grows; U-shape
        (recall test MSEs $3.50,2.60,4.61$ from (b)).
  \end{exercise}
\end{frame}

\begin{frame}{Problem 2 --- solution (a): the decomposition}
  \begin{solutionbox}[Solution P2(a) --- bias--variance decomposition (4 P)]
    At a fixed test point $x_0$ the expected test MSE splits into three non-negative parts:
    \[
      E\big[(y_0-\hat f(x_0))^2\big]
      =\underbrace{\operatorname{Var}(\hat f(x_0))}_{\text{sensitivity to training set}}
      +\underbrace{[\operatorname{Bias}(\hat f(x_0))]^2}_{\text{systematic model error}}
      +\underbrace{\operatorname{Var}(\varepsilon)}_{\text{irreducible noise}}.
    \]
    \begin{itemize}
      \item \textbf{Variance} $\operatorname{Var}(\hat f(x_0))$: how much $\hat f(x_0)$ would
        change if the method were re-fitted on a \emph{different} training set.
      \item \textbf{Squared bias} $[\operatorname{Bias}(\hat f(x_0))]^2$: the systematic error
        from approximating the true $f$ by a (possibly too simple) model.
      \item \textbf{Irreducible error} $\operatorname{Var}(\varepsilon)=\sigma^2$: the noise
        variance of $\varepsilon$, which \emph{no} method can remove.
    \end{itemize}
  \end{solutionbox}
\end{frame}

\begin{frame}{Problem 2 --- solution (b): expected test MSE for A--C}
  \begin{solutionbox}[Solution P2(b) --- compute the test MSE and choose (5 P)]
    Expected test MSE $=$ Bias$^2+\operatorname{Var}+\sigma^2$, with $\sigma^2=1.00$. Adding the
    three components row by row:
    \begin{center}\footnotesize
    \begin{tabular}{lcccc}
      \toprule
       & Bias$^2$ & Var & $\sigma^2$ & Test MSE\\
      \midrule
      A (linear)        & 2.25 & 0.25 & 1.00 & $2.25+0.25+1.00=3.50$\\
      B (moderate)      & 0.16 & 1.44 & 1.00 & $0.16+1.44+1.00=\textbf{2.60}$\\
      C (very flexible) & 0.01 & 3.60 & 1.00 & $0.01+3.60+1.00=4.61$\\
      \bottomrule
    \end{tabular}
    \end{center}
    \textbf{Choose Method B} --- it has the smallest expected test MSE. Its moderate flexibility
    buys a large bias reduction ($2.25\to0.16$) at an acceptable variance cost, whereas the very
    flexible Method C pays more in variance ($+2.16$) than it gains in bias ($-0.15$).
  \end{solutionbox}
\end{frame}

\begin{frame}{Problem 2 --- solution (c): the lower bound}
  \begin{solutionbox}[Solution P2(c) --- smallest achievable test MSE (3 P)]
    Squared bias and variance are both \emph{non-negative}, so for \emph{every} method
    \[
      E\big[(y_0-\hat f(x_0))^2\big]
      =\underbrace{[\operatorname{Bias}(\hat f(x_0))]^2}_{\ge 0}
      +\underbrace{\operatorname{Var}(\hat f(x_0))}_{\ge 0}
      +\sigma^2 \;\ge\; \sigma^2 = 1.00.
    \]
    The floor $1.00$ is attained only in the ideal case of \emph{zero bias and zero variance}.
    Even the true $f$ cannot predict the fresh noise draw $\varepsilon_0$ at $x_0$ --- this
    $\operatorname{Var}(\varepsilon)=\sigma^2$ is the \textbf{irreducible error}.
  \end{solutionbox}
  \begin{takeaway}[Interpretation]
    No modelling effort can push the expected test MSE below $\sigma^2=1.00$.
  \end{takeaway}
\end{frame}

\begin{frame}{Problem 2 --- solution (d): behaviour and the U-shape}
  \begin{solutionbox}[Solution P2(d) --- flexibility and the U-shape (3 P)]
    As the flexibility of a method increases:
    \begin{itemize}
      \item \textbf{Squared bias} decreases (typically monotonically) --- a richer model
        approximates the true $f$ better;
      \item \textbf{Variance} increases --- the fit reacts more to the particular training set;
      \item \textbf{Training MSE} decreases \emph{monotonically} --- a more flexible fit follows
        the training data ever more closely.
    \end{itemize}
    The \textbf{expected test MSE} $=$ Bias$^2+\operatorname{Var}+\sigma^2$ therefore \emph{first
    falls} (bias reduction dominates) and \emph{then rises} (variance increase dominates) ---
    the characteristic \textbf{U-shape}.
    \medskip

    In the table the test MSE runs $3.50 \to 2.60 \to 4.61$ for A $\to$ B $\to$ C: the minimum is
    at the \emph{intermediate} flexibility (B).
  \end{solutionbox}
\end{frame}

% ============================================================
% P3
% ============================================================
\begin{frame}{Problem 3 --- task}
  \begin{exercise}[Problem 3 (12 P) --- KNN classification]
    Seven visitors with scores $X_1$ (pages) and $X_2$ (minutes) and outcome $Y$ (purchase);
    new visitor at $x_0=(4{{,}}2)$ --- classify with Euclidean distance.
    \begin{center}\scriptsize
    \begin{tabular}{lccccccc}
      \toprule
      Obs. & 1 & 2 & 3 & 4 & 5 & 6 & 7\\
      \midrule
      $X_1$ & 4 & 2 & 6 & 4 & 1 & 7 & 5\\
      $X_2$ & 3 & 2 & 2 & 5 & 1 & 4 & 6\\
      $Y$ & Yes & No & No & No & Yes & Yes & No\\
      \bottomrule
    \end{tabular}
    \end{center}
    (a) \textbf{[4 P]} All distances; classification with $K=1$.
    (b) \textbf{[4 P]} Classification with $K=3$ and $\widehat{\Pr}(\text{Yes}\mid x_0)$.
    (c) \textbf{[2 P]} Effect of $K$ on bias/variance; training error of 1-NN.
    (d) \textbf{[2 P]} Bayes classifier: definition; why unattainable; relation to KNN.
  \end{exercise}
\end{frame}

\begin{frame}{Problem 3 --- solution (a): distances and $K=1$}
  \begin{solutionbox}[Solution P3(a) --- Euclidean distances and {$K=1$} (4 P)]
    Distance from $x_0=(4,2)$: $\;d(x_0,x_i)=\sqrt{(X_{1i}-4)^2+(X_{2i}-2)^2}$. Two worked cases:
    \[
      \text{Obs.\ 1:}\ \sqrt{(4-4)^2+(3-2)^2}=\sqrt{1}=1.000,\qquad
      \text{Obs.\ 5:}\ \sqrt{(1-4)^2+(1-2)^2}=\sqrt{10}=3.162.
    \]
    Ranking all seven by squared distance $d^2$ (order-preserving, avoids square roots):
    \begin{center}\scriptsize
    \begin{tabular}{lccccccc}
      \toprule
      Obs. & 1 & 2 & 3 & 4 & 5 & 6 & 7\\
      \midrule
      $d^2$ & 1 & 4 & 4 & 9 & 10 & 13 & 17\\
      $d$ & 1.000 & 2.000 & 2.000 & 3.000 & 3.162 & 3.606 & 4.123\\
      $Y$ & Yes & No & No & No & Yes & Yes & No\\
      \bottomrule
    \end{tabular}
    \end{center}
    \textbf{$K=1$:} the single nearest neighbour is Obs.~1 ($d=1.000$, label Yes)
    $\Rightarrow$ classify the new visitor as \textbf{Yes}.
  \end{solutionbox}
\end{frame}

\begin{frame}{Problem 3 --- solution (b): $K=3$ and the probability}
  \begin{solutionbox}[Solution P3(b) --- majority vote with {$K=3$} (4 P)]
    KNN estimates each class probability by its frequency among the $K$ nearest neighbours:
    \[
      \widehat{\Pr}(Y=j\mid X=x_0)=\frac{1}{K}\sum_{i\in\mathcal N_0}\mathbf{1}(y_i=j),
      \qquad \mathcal N_0=\{\text{$K$ nearest neighbours}\}.
    \]
    From the ranking in part~(a) the three nearest are Obs.~1 ($d=1.000$, Yes),
    Obs.~2 ($d=2.000$, No) and Obs.~3 ($d=2.000$, No). Hence
    \[
      \widehat{\Pr}(\text{Yes}\mid x_0)=\tfrac{1}{3}\approx 0.333,\qquad
      \widehat{\Pr}(\text{No}\mid x_0)=\tfrac{2}{3}\approx 0.667.
    \]
    Majority vote $\Rightarrow$ classify as \textbf{No}. Note that $K=1$ (Yes) and $K=3$ (No)
    \emph{disagree} --- the decision depends on the flexibility of the classifier.
  \end{solutionbox}
\end{frame}

\begin{frame}{Problem 3 --- solution (c): the role of $K$}
  \begin{solutionbox}[Solution P3(c) --- bias/variance and 1-NN training error (2 P)]
    \textbf{Small $K$} (e.g.\ $K=1$): very flexible, jagged decision boundary
    $\Rightarrow$ \emph{low bias, high variance}.\\[0.5mm]
    \textbf{Large $K$}: averaging over many neighbours gives a smooth, almost linear boundary
    $\Rightarrow$ \emph{high bias, low variance}.
    \medskip

    \textbf{Training error of 1-NN $=0$} (ignoring ties): each training point is its \emph{own}
    nearest neighbour, so it is always classified correctly on the training set --- which says
    \emph{nothing} about its test error.
  \end{solutionbox}
  \begin{takeaway}[$K$ is the flexibility knob]
    Small $K$ $=$ flexible / high variance; large $K$ $=$ rigid / high bias --- the same
    bias--variance trade-off as everywhere.
  \end{takeaway}
\end{frame}

\begin{frame}{Problem 3 --- solution (d): the Bayes classifier}
  \begin{solutionbox}[Solution P3(d) --- Bayes classifier and its link to KNN (2 P)]
    \textbf{Definition.} The Bayes classifier assigns each $x$ to the class $j$ with the largest
    \emph{true} conditional probability $\Pr(Y=j\mid X=x)$. It minimises the expected test error
    rate; its error --- the \textbf{Bayes error rate} --- is the lower bound for \emph{all}
    classifiers.
    \medskip

    \textbf{Why it cannot be used directly.} In practice the true conditional distribution of $Y$
    given $X$ is \emph{unknown}, so $\Pr(Y=j\mid X=x)$ cannot be evaluated.
    \medskip

    \textbf{Relation to KNN.} KNN \emph{mimics} the Bayes rule: it \emph{estimates} the
    conditional probabilities locally from the $K$ nearest training observations (the part~(b)
    formula) and then classifies to the largest estimate --- a data-driven approximation of the
    Bayes classifier.
  \end{solutionbox}
\end{frame}

% ============================================================
% P4
% ============================================================
\begin{frame}{Problem 4 --- task}
  \begin{exercise}[Problem 4 (24 P) --- Simple linear regression by hand]
    Radio spots $x$ and weekly sales $y$ (in 1{,}000 EUR) for $n=5$ weeks;
    $\sum x_i=15$; $\sum y_i=55$; $\sum x_i^2=55$; $\sum x_iy_i=195$.
    \begin{center}\scriptsize
    \begin{tabular}{lccccc}
      \toprule
      Week & 1 & 2 & 3 & 4 & 5\\
      \midrule
      $x_i$ & 1 & 2 & 3 & 4 & 5\\
      $y_i$ & 6 & 7 & 11 & 13 & 18\\
      \bottomrule
    \end{tabular}
    \end{center}
    (a) \textbf{[6 P]} $\hat\beta_1$ and $\hat\beta_0$; fitted line; interpretation.
    (b) \textbf{[4 P]} With $\hat y=2+3x$: fitted values and residuals; check $\sum e_i=0$.
    (c) \textbf{[3 P]} With residuals $(1,-1,0,-1,1)$: RSS and RSE; interpretation.
    (d) \textbf{[6 P]} $\operatorname{SE}(\hat\beta_1)$ (use $\hat\beta_1=3$, RSE $=1.155$, $S_{xx}=10$);
        $t$-test of $H_0\colon\beta_1=0$ ($t_{0.975{{,}}3}=3.182$); 95\,\% CI.
    (e) \textbf{[3 P]} TSS and $R^2$ (RSS $=4$); interpretation. \quad
    (f) \textbf{[2 P]} With $\hat y=2+3x$: prediction at $x=6$; one caveat.
  \end{exercise}
\end{frame}

\begin{frame}{Problem 4 --- solution (a): least-squares estimates}
  \begin{solutionbox}[Solution P4(a) --- slope $\hat\beta_1$ and intercept $\hat\beta_0$ (6 P)]
    \textbf{Step 1 --- means:}\quad
    $\bar x=\tfrac{\sum x_i}{n}=\tfrac{15}{5}=3$,\qquad
    $\bar y=\tfrac{\sum y_i}{n}=\tfrac{55}{5}=11$.

    \textbf{Step 2 --- corrected cross-products:}
    \begin{align*}
      S_{xy}&=\sum x_iy_i-n\bar x\bar y=195-5\cdot 3\cdot 11=195-165=30,\\
      S_{xx}&=\sum x_i^2-n\bar x^2=55-5\cdot 3^2=55-45=10.
    \end{align*}

    \textbf{Step 3 --- slope and intercept:}
    \[
      \hat\beta_1=\frac{S_{xy}}{S_{xx}}=\frac{30}{10}=3.000,\qquad
      \hat\beta_0=\bar y-\hat\beta_1\bar x=11-3\cdot 3=2.000,
    \]
    giving the fitted line $\hat y=2.000+3.000\,x$.

    \textbf{Interpretation.} Slope: each extra weekly spot adds an expected $3{,}000$ EUR in
    sales. Intercept: expected sales of $2{,}000$ EUR at zero spots --- a mild extrapolation
    ($x=0$ unobserved; $x$ ranges 1--5).
  \end{solutionbox}
\end{frame}

\begin{frame}{Problem 4 --- solution (b): fitted values and residuals}
  \begin{solutionbox}[Solution P4(b) --- fitted values $\hat y_i$ and residuals $e_i$ (4 P)]
    With $\hat y_i=2+3x_i$ and $e_i=y_i-\hat y_i$; e.g.\ $\hat y_1=2+3\cdot 1=5$ so $e_1=6-5=1$:
    \begin{center}\footnotesize
    \begin{tabular}{lccccc}
      \toprule
      $i$ & 1 & 2 & 3 & 4 & 5\\
      \midrule
      $x_i$ & 1 & 2 & 3 & 4 & 5\\
      $y_i$ & 6 & 7 & 11 & 13 & 18\\
      $\hat y_i=2+3x_i$ & 5 & 8 & 11 & 14 & 17\\
      $e_i=y_i-\hat y_i$ & $1$ & $-1$ & $0$ & $-1$ & $1$\\
      \bottomrule
    \end{tabular}
    \end{center}
    \textbf{Check:} $\sum_{i=1}^{5} e_i = 1-1+0-1+1 = 0$ \checkmark
    \smallskip

    This is no coincidence: for any least-squares fit \emph{with an intercept}, the residuals
    always sum to zero (one of the normal-equation identities).
  \end{solutionbox}
\end{frame}

\begin{frame}{Problem 4 --- solution (c): RSS and RSE}
  \begin{solutionbox}[Solution P4(c) --- residual sum of squares and RSE (3 P)]
    \textbf{Residual sum of squares} (sum of squared residuals from part~(b)):
    \[
      \operatorname{RSS}=\sum_{i=1}^{5} e_i^2 = 1^2+(-1)^2+0^2+(-1)^2+1^2 = 1+1+0+1+1 = 4.000.
    \]
    \textbf{Residual standard error} (with $n-2=3$ degrees of freedom):
    \[
      \operatorname{RSE}=\sqrt{\frac{\operatorname{RSS}}{\,n-2\,}}=\sqrt{\frac{4}{3}}
      =\sqrt{1.333}=1.155 \ \text{(thousand EUR)}.
    \]
    \textbf{Interpretation.} Even if the true regression line were known, observed weekly sales
    would still deviate from it by about $1.155$ thousand EUR ($\approx 1{,}155$ EUR) on average.
    RSE is the estimate $\hat\sigma$ of $\sigma$, the standard deviation of the error $\varepsilon$.
  \end{solutionbox}
\end{frame}

\begin{frame}{Problem 4 --- solution (d.1): standard error and $t$-test}
  \begin{solutionbox}[Solution P4(d) --- standard error and hypothesis test (part 1 of 2)]
    \textbf{Hypotheses:} $H_0\colon\beta_1=0$ vs.\ $H_1\colon\beta_1\ne 0$, at level $\alpha=0.05$.
    \medskip

    \textbf{Standard error of the slope:}
    \[
      \operatorname{SE}(\hat\beta_1)=\frac{\operatorname{RSE}}{\sqrt{S_{xx}}}
      =\frac{1.1547}{\sqrt{10}}=\frac{1.1547}{3.1623}=0.365.
    \]
    \textbf{Test statistic} (df $=n-2=3$):
    \[
      t=\frac{\hat\beta_1-0}{\operatorname{SE}(\hat\beta_1)}=\frac{3.000}{0.3651}=8.216.
    \]
    \textbf{Decision:} since $|t|=8.216>3.182=t_{0.975,\,3}$, we \emph{reject} $H_0$ at the
    $5\,\%$ level --- the slope is significantly different from zero.
  \end{solutionbox}
\end{frame}

\begin{frame}{Problem 4 --- solution (d.2): confidence interval}
  \begin{solutionbox}[Solution P4(d) --- 95\,\% confidence interval for $\beta_1$ (part 2 of 2)]
    Using $\hat\beta_1=3.000$, $\operatorname{SE}(\hat\beta_1)=0.3651$ and $t_{0.975,\,3}=3.182$:
    \begin{align*}
      \hat\beta_1 \pm t_{0.975,\,3}\cdot\operatorname{SE}(\hat\beta_1)
      &= 3.000 \pm 3.182\cdot 0.3651\\
      &= 3.000 \pm 1.162\\
      &= [\,1.838,\ 4.162\,].
    \end{align*}
    Zero is \emph{not} contained in the interval --- consistent with rejecting $H_0$ in
    part~(d.1). \textbf{Conclusion:} there is a statistically significant \emph{positive} linear
    association between the number of radio spots and weekly sales.
  \end{solutionbox}
  \begin{takeaway}[Test $\Leftrightarrow$ CI]
    A two-sided $5\,\%$ test rejects $H_0\colon\beta_1=0$ exactly when $0$ falls outside the
    $95\,\%$ confidence interval.
  \end{takeaway}
\end{frame}

\begin{frame}{Problem 4 --- solution (e): TSS and $R^2$}
  \begin{solutionbox}[Solution P4(e) --- total sum of squares and $R^2$ (3 P)]
    \textbf{Total sum of squares} (deviations of $y_i$ from $\bar y=11$):
    \[
      \operatorname{TSS}=\sum_{i=1}^{5}(y_i-\bar y)^2
      =(-5)^2+(-4)^2+0^2+2^2+7^2 = 25+16+0+4+49 = 94.000.
    \]
    \textbf{Coefficient of determination} (using $\operatorname{RSS}=4.000$):
    \[
      R^2=1-\frac{\operatorname{RSS}}{\operatorname{TSS}}=1-\frac{4}{94}=\frac{90}{94}=0.957.
    \]
    \textbf{Interpretation.} About $95.7\,\%$ of the variance of weekly sales is explained by the
    linear regression on the number of radio spots. (In simple regression $R^2=r^2$, so the
    correlation is $r=\sqrt{0.957}\approx 0.979$.)
  \end{solutionbox}
\end{frame}

\begin{frame}{Problem 4 --- solution (f): prediction at $x=6$}
  \begin{solutionbox}[Solution P4(f) --- prediction and its caveat (2 P)]
    \textbf{Point prediction} from $\hat y=2+3x$ at $x=6$:
    \[
      \hat y(6)=2+3\cdot 6 = 2+18 = 20.000 \ \text{thousand EUR} = 20{,}000\ \text{EUR}.
    \]
    \textbf{Caveat.} $x=6$ lies \emph{outside} the observed range $x\in[1,5]$ --- this is an
    \textbf{extrapolation}. The linear relationship has not been observed there and may not hold;
    moreover the sample $n=5$ is very small, so the estimates are themselves uncertain.
  \end{solutionbox}
  \begin{takeaway}[Extrapolation warning]
    A fitted line is only trustworthy \emph{within} the range of the observed predictor values.
  \end{takeaway}
\end{frame}

% ============================================================
% P5
% ============================================================
\begin{frame}{Problem 5 --- task}
  \begin{exercise}[Problem 5 (24 P) --- Multiple regression output]
    $n=60$ employees; \texttt{salary} (1{,}000 EUR/year); \texttt{experience} (years);
    \texttt{female} (dummy). OLS fit of
    $\texttt{salary}=\beta_0+\beta_1\texttt{experience}+\beta_2\texttt{female}
      +\beta_3(\texttt{experience}\!\times\!\texttt{female})+\varepsilon$:
    \begin{center}\scriptsize
    \begin{tabular}{lrrrr}
      \toprule
      & coef & std err & $t$ & $P>|t|$\\
      \midrule
      Intercept & 46.2000 & 1.100 & 42.000 & $<$0.0001\\
      experience & 2.3500 & 0.250 & 9.400 & $<$0.0001\\
      female & $-3.1000$ & 1.550 & $-2.000$ & 0.0504\\
      experience:female & $-0.8500$ & 0.350 & $-2.429$ & 0.0184\\
      \bottomrule
    \end{tabular}\\[1mm]
    $R^2=0.712$;\ adj.\ $R^2=0.697$;\ $F=46.15$ on 3 and 56 df; Prob($F$) $<0.0001$
    \end{center}
    (a) \textbf{[6 P]} Interpret all coefficients. (b) \textbf{[4 P]} Significance at 5\,\%; drop \texttt{female}?
    (c) \textbf{[4 P]} Fitted equation per group. (d) \textbf{[4 P]} Predictions at 10 years; gap.
    (e) \textbf{[3 P]} $R^2$ and $F$-test. (f) \textbf{[3 P]} U-shaped residual plot; VIF $=12$.
  \end{exercise}
\end{frame}

\begin{frame}{Problem 5 --- solution (a.1): intercept and experience}
  \begin{solutionbox}[Solution P5(a) --- baseline coefficients (part 1 of 2)]
    Salary is measured in $1{,}000$ EUR per year; the \emph{baseline} group is men
    (\texttt{female}$=0$).
    \medskip

    \textbf{Intercept $=46.20$.} Expected annual salary of a \emph{man} with $0$ years of
    experience: $46{,}200$ EUR.
    \medskip

    \textbf{\texttt{experience} $=2.35$.} For \emph{men}, each additional year of experience
    raises the expected salary by $2{,}350$ EUR. Because the model contains an interaction, this
    slope refers to the baseline group (\texttt{female}$=0$) \emph{only} --- women have a
    different slope (see part~a.2).
  \end{solutionbox}
  \begin{takeaway}[Reading a dummy--interaction model]
    With \texttt{female} and \texttt{experience:female} present, \texttt{Intercept} and
    \texttt{experience} describe the \emph{men} only.
  \end{takeaway}
\end{frame}

\begin{frame}{Problem 5 --- solution (a.2): the dummy and the interaction}
  \begin{solutionbox}[Solution P5(a) --- gender difference and its slope (part 2 of 2)]
    \textbf{\texttt{female} $=-3.10$.} At $0$ years of experience, women are expected to earn
    $3{,}100$ EUR \emph{less} than men --- the difference between the two group \emph{intercepts}.
    \medskip

    \textbf{\texttt{experience:female} $=-0.85$.} The experience slope for women is $0.85$
    \emph{lower} than for men:
    \[
      \text{women's slope}=2.35-0.85=1.50\ \text{thousand EUR per year}.
    \]
    Equivalently, the expected salary gap between men and women \emph{widens} by $850$ EUR with
    every additional year of experience.
  \end{solutionbox}
\end{frame}

\begin{frame}{Problem 5 --- solution (b): significance and the hierarchy principle}
  \begin{solutionbox}[Solution P5(b) --- significance and the hierarchy principle (4 P)]
    At $\alpha=0.05$ compare each $p$-value with $0.05$:
    \begin{itemize}
      \item \texttt{Intercept}, \texttt{experience}: $p<0.0001$ $\Rightarrow$ \emph{highly significant}.
      \item \texttt{experience:female}: $p=0.0184<0.05$ $\Rightarrow$ \emph{significant}.
        Consistency check: $t=\hat\beta_3/\operatorname{SE}(\hat\beta_3)=-0.85/0.35=-2.429$
        \checkmark\ (matches the output).
      \item \texttt{female}: $p=0.0504>0.05$ $\Rightarrow$ \emph{not} significant at exactly the
        $5\,\%$ level (only marginally).
    \end{itemize}
    \textbf{Should \texttt{female} be removed? No.} By the \textbf{hierarchy principle}, if an
    interaction (\texttt{experience}$\times$\texttt{female}) is kept, its main effects must be
    kept too --- otherwise the interaction changes meaning. Besides, $p=0.0504$ is borderline
    evidence, not evidence of \emph{no} effect.
  \end{solutionbox}
\end{frame}

\begin{frame}{Problem 5 --- solution (c): fitted equation per group}
  \begin{solutionbox}[Solution P5(c) --- separate lines for men and women (4 P)]
    Substitute \texttt{female}$=0$ and \texttt{female}$=1$ into
    $\widehat{\texttt{salary}}=46.20+2.35\,\texttt{exp}-3.10\,\texttt{female}
      -0.85\,(\texttt{exp}\times\texttt{female})$:
    \begin{align*}
      \text{Men } (\texttt{female}=0):\quad
        & \widehat{\texttt{salary}} = 46.20 + 2.35\,\texttt{experience}\\
      \text{Women } (\texttt{female}=1):\quad
        & \widehat{\texttt{salary}} = (46.20-3.10) + (2.35-0.85)\,\texttt{experience}\\
        & \phantom{\widehat{\texttt{salary}}} = 43.10 + 1.50\,\texttt{experience}
    \end{align*}
    The two groups get \emph{different intercepts and different slopes} --- exactly what including
    the interaction term allows.
  \end{solutionbox}
\end{frame}

\begin{frame}{Problem 5 --- solution (d): predictions and the gap at 10 years}
  \begin{solutionbox}[Solution P5(d) --- predicted salaries at 10 years (4 P)]
    Insert \texttt{experience}$=10$ into the group equations from part~(c):
    \begin{align*}
      \text{Man:}\quad & 46.20 + 2.35\cdot 10 = 46.20+23.50 = 69.70 \Rightarrow 69{,}700\ \text{EUR}\\
      \text{Woman:}\quad & 43.10 + 1.50\cdot 10 = 43.10+15.00 = 58.10 \Rightarrow 58{,}100\ \text{EUR}
    \end{align*}
    \textbf{Estimated gap} (woman $-$ man): $58.10-69.70=-11.60$ thousand EUR --- women earn
    $11{,}600$ EUR less. In terms of the coefficients the gap at experience $=10$ is
    \[
      \hat\beta_2+10\,\hat\beta_3=-3.10-8.50=-11.60\ \text{thousand EUR}.
    \]
  \end{solutionbox}
\end{frame}

\begin{frame}{Problem 5 --- solution (e): $R^2$ and the $F$-test}
  \begin{solutionbox}[Solution P5(e) --- overall fit and joint significance (3 P)]
    \textbf{$R^2=0.712$.} The three model terms (\texttt{experience}, \texttt{female} and their
    interaction) jointly explain $71.2\,\%$ of the variance of salaries in the sample.
    \medskip

    \textbf{$F$-test --- is the model useful at all?}
    \[
      H_0\colon \beta_1=\beta_2=\beta_3=0 \quad\text{vs.}\quad
      H_1\colon \text{at least one } \beta_j\neq 0.
    \]
    The output gives $F=46.15$ on $3$ and $56$ degrees of freedom with $p<0.0001$
    $\Rightarrow$ \emph{reject} $H_0$: at least one predictor is related to salary, so the model
    is useful overall.
  \end{solutionbox}
\end{frame}

\begin{frame}{Problem 5 --- solution (f.1): non-linearity}
  \begin{solutionbox}[Solution P5(f)(i) --- U-shaped residual plot (part 1 of 2)]
    A residual-vs-fitted plot with a clear \textbf{U-shape} (positive residuals at small and
    large fitted values, negative in between) signals \textbf{non-linearity}: the straight-line
    model misses systematic \emph{curvature}, so the residuals are not patternless as they should be.
    \medskip

    \textbf{Remedies.}
    \begin{itemize}
      \item Add a non-linear term, e.g.\ $\texttt{experience}^2$ (polynomial regression);
      \item Transform variables, e.g.\ model $\log(\texttt{salary})$ or use $\sqrt{X}$.
    \end{itemize}
  \end{solutionbox}
  \begin{takeaway}[Diagnostic rule]
    A good linear fit shows a \emph{structureless} residual cloud; any systematic shape points to
    a violated assumption.
  \end{takeaway}
\end{frame}

\begin{frame}{Problem 5 --- solution (f.2): multicollinearity (VIF)}
  \begin{solutionbox}[Solution P5(f)(ii) --- what {$\operatorname{VIF}=12$} means (part 2 of 2)]
    \textbf{Definition.} $\operatorname{VIF}_j=\dfrac{1}{1-R^2_{X_j\mid X_{-j}}}$, where
    $R^2_{X_j\mid X_{-j}}$ comes from regressing $X_j$ on all \emph{other} predictors. Solving
    $\operatorname{VIF}_{\texttt{age}}=12$:
    \[
      R^2_{\texttt{age}\mid \text{rest}}=1-\frac{1}{12}=0.917.
    \]
    So \texttt{age} is almost a linear combination of the other predictors --- strong
    \textbf{multicollinearity} (common rules of thumb flag VIF $>5$ or $>10$).
    \medskip

    \textbf{Consequence.} $\operatorname{SE}(\hat\beta_{\texttt{age}})$ is inflated by a factor
    $\sqrt{12}\approx 3.5$: the coefficient is estimated very imprecisely, its sign may flip
    between samples, and its $t$-test loses power --- even though the model's overall fit and
    predictions may be unaffected.
  \end{solutionbox}
\end{frame}

\end{document}
